% Reconstructed from the compiled lecture-note PDF.
\chapter{The Pontryagin system for Reinhardt's problem}
\section{State and costate variables}

The state is (g, X) ∈SL2(R) × sl2(R).


We represent the two costates by matrices


\begin{displaytext}
Λ1, Λ2 ∈sl2(R)
\end{displaytext}

\begin{displaytext}
through the trace pairing. The costate Λ1 corresponds to the group variable g, while Λ2 corresponds \
to X.
\end{displaytext}

\begin{displaytext}
The dynamics are \
Zu g′ = gX, X′ = [P, X], P = ⟨Zu, X⟩. \
The running cost is \
l(X) = −3 ⟨J, X⟩. \
2
\end{displaytext}

\section{The Hamiltonian}

The group part contributes ⟨Λ1, X⟩.


\begin{displaytext}
The X-state part contributes \
⟨Λ2, [P, X]⟩= −⟨[Λ2, X], Zu⟩ . \
⟨X, Zu⟩ \
It is convenient to define the reduced costate
\end{displaytext}

\begin{displaytext}
ΛR = [Λ2, X].
\end{displaytext}

Then the full Hamiltonian is


\begin{displaytext}
Zu⟩ H(Λ1, ΛR, X; u) = Λ1 −3 X −⟨ΛR, 2λcostJ, ⟨X, Zu⟩.
\end{displaytext}

\begin{displaytext}
For normal extremals we take λcost = −1, so the first term becomes
\end{displaytext}

3 Λ1 + 2J, X .


The source often keeps λcost symbolic until normality is established.


\begin{insightbox}{Point requiring care}
\end{insightbox}

\begin{displaytext}
The Pontryagin convention maximizes H and uses λcost ≤0. Other textbooks minimize the \
Hamiltonian or use a positive cost multiplier. Signs must be translated before comparing \
formulas.
\end{displaytext}

\section{The state–costate equations}

Let Zu P = ⟨Zu, X⟩.


The complete equations are


\begin{displaytext}
g′ = gX, (9.1) \
X′ = [P, X], (9.2) \
Λ′1 = [Λ1, X], (9.3) \
3 \
Λ′R = [P, ΛR] −⟨ΛR, P⟩[P, X] + −Λ1 + 2λcostJ, X . (9.4)
\end{displaytext}

The control u(t) maximizes H over UT almost everywhere, and


H = 0


along a free-time extremal.


The equation for Λ1 is another Lax equation. Its solution is explicit:


\begin{displaytext}
Λ1(t) = g(t)−1Λ1(0)g(t).
\end{displaytext}

\begin{displaytext}
Thus \
det Λ1(t) = det Λ1(0)
\end{displaytext}

is conserved.


\section{Why ΛR is the right reduced variable}

\begin{displaytext}
By construction, \
⟨ΛR, X⟩= ⟨[Λ2, X], X⟩= 0.
\end{displaytext}

Hence ΛR lies in the two-dimensional plane orthogonal to X under the trace form. This is exactly the cotangent space to the determinant-one orbit at X.


The control-dependent part of H depends only on (X, ΛR):


\begin{displaytext}
Zu⟩ H2(X, ΛR; u) = −⟨ΛR, ⟨X, Zu⟩.
\end{displaytext}

\begin{displaytext}
The variables g and Λ1 matter for endpoint conditions and the inhomogeneous term in Equation (9.4), \
but the switching geometry is controlled by (X, ΛR).
\end{displaytext}

\section{Hamiltonian maximizers are faces}

\begin{displaytext}
Fix (X, ΛR). Along a segment in the control simplex, Zu is affine in u, so H2 is a ratio of affine \
functions. As explained in Chapter 8, the set of maximizers is a face. Thus the cotangent space is \
partitioned into seven strata indexed by nonempty subsets I ⊂\{0, 1, 2\}:
\end{displaytext}

\begin{displaytext}
arg max H2 = (UT )I, \
u∈UT
\end{displaytext}

where (UT )I is the face spanned by the vertices in I.


The three open vertex regions are separated by three codimension-one walls. The walls meet at the full singular set where the Hamiltonian is independent of u.


\section{Switching functions}

Let Zj Pj(X) = ⟨Zj, X⟩


for the three vertex controls. The difference between the control-dependent Hamiltonians at vertices i and j is, up to a harmless sign convention,


\begin{displaytext}
χij(X, ΛR) = ⟨ΛR, Pj(X) −Pi(X)⟩.
\end{displaytext}

A switch between i and j occurs on the wall


\begin{displaytext}
χij = 0.
\end{displaytext}

Away from multiple zeros, the first switching time is a smooth function of the initial state and costate.


\section{Transversality and periodicity}

The geometric endpoint conditions are rotational:


\begin{displaytext}
g(tf) = R, X(tf) = R−1X(0)R.
\end{displaytext}

The costates satisfy matching conditions of the same form:


\begin{displaytext}
Λ1(tf) = R−1Λ1(0)R, ΛR(tf) = R−1ΛR(0)R.
\end{displaytext}

Extending by symmetry yields


\begin{displaytext}
X(t + 3tf) = X(t), Λ1(t + 3tf) = Λ1(t), ΛR(t + 3tf) = ΛR(t).
\end{displaytext}

The group variable has period 6tf because R6 = I.


\begin{insightbox}{Geometric intuition}
\end{insightbox}

The final contradiction will use only one global feature of the boundary conditions: a minimizer extends to a periodic Pontryagin trajectory. The local chattering analysis will show that every trajectory entering the singular locus lies on a stable curve that cannot close periodically.


\section{Left invariance}

\begin{displaytext}
The equations and cost depend on g only through g′ = gX, and replacing g(t) by hg(t) with constant \
h ∈SL2(R) changes neither X nor the cost. This is left invariance. It explains why the group \
costate satisfies a conjugation equation and why much of the dynamical analysis can be performed \
solely in the Lie algebra variables.
\end{displaytext}

\begin{displaytext}
For the proof of Mahler’s First conjecture, one often temporarily ignores the endpoint condition on \
g, analyzes (X, Λ1, ΛR), and restores g(tf) = R at the end. This separation is legitimate because g \
is reconstructed by integrating g′ = gX once X is known.
\end{displaytext}

\section{Normality near the singular locus}

\begin{displaytext}
At a point where \
ΛR = 0, Λ′R = 0,
\end{displaytext}

the Hamiltonian equation and Equation (9.4) force


\begin{displaytext}
3 \
Λ1 = 2λcostJ.
\end{displaytext}

\begin{displaytext}
If λcost = 0, then both costates vanish, contradicting the nontriviality condition in the maximum \
principle. Therefore \
λcost ̸= 0.
\end{displaytext}

\begin{displaytext}
Any extremal that reaches the full singular locus is normal, and after scaling \
λcost = −1, Λ1 = −3 2J.
\end{displaytext}

These fixed parameter values are used in the blowup analysis.


\section{Constant vertex segments}

\begin{displaytext}
When u = ej is constant, P = Pj is constant because ⟨X, Zj⟩is conserved. Hence X and g have \
explicit solutions:
\end{displaytext}

X(t) = etP X0e−tP ,


g(t) = et(X0+P)e−tP


\begin{displaytext}
when g(0) = I. The costate Λ1 follows by conjugation, and ΛR is obtained by elementary quadratures \
of matrix exponentials.
\end{displaytext}

These explicit formulas make the Poincaré maps and switching functions computable. The eventual computer-assisted part of the proof depends heavily on the fact that constant-control solutions are analytic and explicit.


\begin{insightbox}{Chapter summary}
\end{insightbox}

\begin{displaytext}
The maximum principle lifts the Reinhardt state to the matrix system (g, X, Λ1, ΛR). The \
Hamiltonian is \
Zu⟩ H = Λ1 −3 X −⟨ΛR, 2λJ, ⟨X, Zu⟩,
\end{displaytext}

and the control maximizers are faces of the triangle. Rotational transversality makes every minimizing extremal periodic. The full singular locus is necessarily normal and has fixed costate parameters.

